% =============================================================================
%  APÊNDICE: SOLUÇÕES DOS EXERCÍCIOS
%  Capítulos 1 a 8
% =============================================================================
\chapter{Soluções dos Exercícios}\label{ap:solucoes}

Este apêndice apresenta soluções comentadas para os exercícios dos Capítulos~1
a~8, organizadas por capítulo e seção. As soluções priorizam a didática: cada
resposta explica o raciocínio, não apenas o resultado final.

% =============================================================================
%  CAPÍTULO 1: FUNDAMENTOS MATEMÁTICOS E ESTATÍSTICOS
% =============================================================================
\section*{Capítulo 1 -- Fundamentos Matemáticos e Estatísticos}
\addcontentsline{toc}{section}{Capítulo 1}

\subsection*{Exercício 1.1 (Nível 0) -- Tabela-Verdade}
Construa a tabela-verdade de $(p \land q) \to \neg r$.

\textbf{Solução:} A tabela possui $2^3 = 8$ linhas. A condicional
$A \to B$ é falsa apenas quando $A$ é verdadeiro e $B$ é falso.
Neste caso, $A = p \land q$ e $B = \neg r$:

\[
\begin{array}{ccc|c|c|c}
p & q & r & p \land q & \neg r & (p \land q) \to \neg r \\ \hline
V & V & V & V & F & F \\
V & V & F & V & V & V \\
V & F & V & F & F & V \\
V & F & F & F & V & V \\
F & V & V & F & F & V \\
F & V & F & F & V & V \\
F & F & V & F & F & V \\
F & F & F & F & V & V
\end{array}
\]

A única linha que torna a expressão falsa é $p=V$, $q=V$, $r=V$.

\subsection*{Exercício 1.2 (Nível Básico) -- Lógica de Predicados}
Traduza para lógica de predicados: ``todo agente que falha no teste de
confiança é redirecionado para o modo shadow''.

\textbf{Solução:} Seja $A(x)$ o predicado ``$x$ é um agente'', $F(x)$
``$x$ falha no teste de confiança'' e $S(x)$ ``$x$ é redirecionado para
o modo shadow''. A tradução é:

\[
\forall x\,[(A(x) \land F(x)) \to S(x)]
\]

Lê-se: ``para todo $x$, se $x$ é um agente e $x$ falha no teste de
confiança, então $x$ é redirecionado para o modo shadow''.

\subsection*{Exercício 1.3 (Nível Intermediário) -- Prova por Indução}
Prove que o número de linhas de uma tabela-verdade com $n$ proposições é $2^n$.

\textbf{Solução:} Demonstração por indução em $n$.

\textbf{Base ($n=1$):} Com uma proposição $p$, há duas possibilidades
($V$ ou $F$). Logo, $2^1 = 2$ linhas. $\checkmark$

\textbf{Passo:} Suponha que para $n$ proposições haja $2^n$ linhas
(hipótese de indução). Para $n+1$ proposições, fixamos a $(n+1)$-ésima
proposição em $V$ (o que gera $2^n$ linhas para as $n$ primeiras) e
depois em $F$ (outras $2^n$ linhas). Total: $2^n + 2^n = 2 \cdot 2^n = 2^{n+1}$.
Portanto, a fórmula vale para todo $n \in \mathbb{N}$.

\subsection*{Exercício 1.4 (Nível Avançado) -- Tabela-Verdade em Python}
Implemente em Python uma função que recebe uma expressão booleana como
string e retorna sua tabela-verdade.

\textbf{Solução:}
\begin{lstlisting}[style=pythonstyle]
import itertools

def tabela_verdade(expr: str, vars: list[str]) -> list[dict]:
    resultados = []
    for valores in itertools.product([True, False], repeat=len(vars)):
        env = dict(zip(vars, valores))
        resultado = eval(expr, {}, env)
        linha = {v: env[v] for v in vars}
        linha[expr] = resultado
        resultados.append(linha)
    return resultados

# Teste com a expressao do Trust Engine: (a and t) or not s
expr = "(a and t) or not s"
vars = ["a", "t", "s"]
for linha in tabela_verdade(expr, vars):
    print(linha)
\end{lstlisting}

\subsection*{Exercício 1.5 (Nível Básico) -- Derivada da Sigmoide}
Calcule a derivada de $\sigma(x) = \dfrac{1}{1 + e^{-x}}$ e mostre que
$\sigma'(x) = \sigma(x)(1 - \sigma(x))$.

\textbf{Solução:} Usando a regra da cadeia:

\[
\sigma(x) = (1 + e^{-x})^{-1}
\]
\[
\sigma'(x) = -1 \cdot (1 + e^{-x})^{-2} \cdot (-e^{-x})
= \frac{e^{-x}}{(1 + e^{-x})^2}
\]

Fatorando: $\sigma'(x) = \frac{1}{1 + e^{-x}} \cdot \frac{e^{-x}}{1 + e^{-x}}$.
Como $\sigma(x) = \frac{1}{1 + e^{-x}}$ e $1 - \sigma(x) = \frac{e^{-x}}{1 + e^{-x}}$,
temos $\sigma'(x) = \sigma(x)(1 - \sigma(x))$.

\subsection*{Exercício 1.6 (Nível Intermediário) -- Intervalo de Confiança}
Dada uma amostra de 30 scores do TrustScorer com média 0,74 e desvio
padrão 0,12, construa um IC de 95\% para a média populacional.

\textbf{Solução:} Para $n=30$, usamos a distribuição $t$ com $29$ graus de
liberdade. O quantil $t_{0,025; 29} \approx 2,045$. O erro padrão é
$EP = s / \sqrt{n} = 0,12 / \sqrt{30} \approx 0,0219$. Logo:

\[
IC_{95\%} = \bar{x} \pm t \cdot EP = 0,74 \pm 2,045 \times 0,0219
= 0,74 \pm 0,0448 = [0,6952;\; 0,7848]
\]

Interpretação: há 95\% de confiança de que a média populacional dos scores
do TrustScorer está entre 0,695 e 0,785.

% =============================================================================
%  CAPÍTULO 2: IA E AGENTES INTELIGENTES
% =============================================================================
\section*{Capítulo 2 -- Inteligência Artificial e Agentes Inteligentes}
\addcontentsline{toc}{section}{Capítulo 2}

\subsection*{Exercício 2.1 (Nível 0) -- Problema de Turing}
Explique por que a pergunta ``esta máquina é inteligente?'' é considerada
filosoficamente problemática segundo o argumento de Turing.

\textbf{Solução:} Turing (1950) propôs substituir a pergunta ``máquinas
pensam?'' por um teste comportamental (o Jogo da Imitação). O problema
é: (a) não há definição consensual de ``inteligência''; (b) o Teste de
Turing mede \emph{simulação} de comportamento humano, não consciência;
(c) a pergunta é autorreferente --- quem julga a inteligência também
usa seu próprio critério subjetivo. Por isso, Turing argumentou que é
mais produtivo perguntar ``a máquina pode passar no teste?'' do que
``a máquina é inteligente?''.

\subsection*{Exercício 2.2 (Nível Básico) -- Agente Reativo Simples}
Implemente em Python um agente reativo simples para um termostato.

\textbf{Solução:}
\begin{lstlisting}[style=pythonstyle]
class Termostato:
    def __init__(self, temp_alvo: float = 22.0):
        self.temp_alvo = temp_alvo
        self.ligado = False

    def perceber(self, temp_atual: float) -> str:
        if temp_atual < self.temp_alvo - 1.0:
            return "frio"
        elif temp_atual > self.temp_alvo + 1.0:
            return "quente"
        return "ok"

    def agir(self, percepcao: str) -> str:
        if percepcao == "frio":
            self.ligado = True
            return "ligar_aquecimento"
        elif percepcao == "quente":
            self.ligado = False
            return "desligar_aquecimento"
        return "manter"

    def ciclo(self, temp_atual: float) -> str:
        p = self.perceber(temp_atual)
        return self.agir(p)

t = Termostato()
for temp in [18, 21, 23, 19, 22]:
    acao = t.ciclo(temp)
    print(f"Temp={temp}C -> {acao}")
\end{lstlisting}

\subsection*{Exercício 2.3 (Nível Intermediário) -- Atenção Scaled Dot-Product}
Implemente a função de atenção por produto escalar do zero usando NumPy.

\textbf{Solução:}
\begin{lstlisting}[style=pythonstyle]
import numpy as np

def scaled_dot_product_attention(Q, K, V):
    """Atencao scaled dot-product.\n    Q, K, V: matrizes (n x d_k), (m x d_k), (m x d_v)"""
    d_k = Q.shape[-1]
    scores = Q @ K.T / np.sqrt(d_k)  # (n x m)
    pesos = np.exp(scores) / np.sum(np.exp(scores), axis=-1, keepdims=True)
    return pesos @ V  # (n x d_v)

# Teste com Q, K, V 4x8
np.random.seed(42)
Q = np.random.randn(4, 8)
K = np.random.randn(4, 8)
V = np.random.randn(4, 8)
saida = scaled_dot_product_attention(Q, K, V)
print(f"Saida shape: {saida.shape}")  # (4, 8)
print(f"Saida[0]: {saida[0][:4]}...")
\end{lstlisting}

\subsection*{Exercício 2.4 (Nível Avançado) -- Z3 para Behavioral Gate}
Usando Z3, verifique a consistência das regras do Behavioral Gate.

\textbf{Solução:}
\begin{lstlisting}[style=pythonstyle]
from z3 import *

trust = Real('trust')
acao_permitida = Bool('acao_permitida')
acao_bloqueada = Bool('acao_bloqueada')

regras = [
    Implies(trust >= 0.7, acao_permitida),
    Implies(trust < 0.3, acao_bloqueada),
    Not(And(acao_permitida, acao_bloqueada))
]

solver = Solver()
solver.add(regras)

if solver.check() == sat:
    print("Regras consistentes. Modelo:")
    print(solver.model())
else:
    print("INCONSISTENCIA: conflito entre regras!")
\end{lstlisting}

O Z3 confirma que as regras são consistentes: nunca há um estado em que
ambas as ações (permitida e bloqueada) sejam verdadeiras simultaneamente.

% =============================================================================
%  CAPÍTULO 3: OPENCODE ARCHITECTURE
% =============================================================================
\section*{Capítulo 3 -- OpenCode Arquitetura}
\addcontentsline{toc}{section}{Capítulo 3}

\subsection*{Exercício 3.1 (Nível 0) -- Instalação}
Instale o \opencodeeco{} e execute \codigo{python scripts/validate\_installation.py}.

\textbf{Solução:} A instalação segue os passos:

\begin{enumerate}
\item Clone o repositório: \codigo{git clone https://github.com/MarceloClaro/OpenCode_Ecosystem.git}
\item Acesse o diretório: \codigo{cd opencode}
\item Instale dependências: \codigo{pip install -e .}
\item Execute a validação:
\begin{lstlisting}[style=pythonstyle]
python scripts/validate_installation.py
\end{lstlisting}
\end{enumerate}

O script verifica: (a) presença do Python 3.10+, (b) dependências
instaladas, (c) estrutura de diretórios, (d) permissões de escrita.
Uma saída esperada é: \codigo{Instalacao validada com sucesso! 312/312 CTs OK}.

\subsection*{Exercício 3.2 (Nível Intermediário) -- SDD+TDD}
Implemente a SPEC-039 completa seguindo o ciclo SDD+TDD.

\textbf{Solução (esqueleto):} O ciclo SDD+TDD segue quatro passos:

\begin{enumerate}
\item \textbf{SDD -- Especificação:} Defina a SPEC com requisitos funcionais
(RF01--RF05), assinaturas e invariantes.
\item \textbf{TDD -- Teste:} Antes de implementar, escreva um CT que capture
o comportamento esperado:
\begin{lstlisting}[style=pythonstyle]
def test_sumarizador_retorna_resumo():
    s = AutoSummarizer(max_length=100)
    resultado = s.summarize("Texto longo...")
    assert len(resultado) <= 100
    assert isinstance(resultado, str)
\end{lstlisting}
\item \textbf{Implementação:} Codifique a função mínima para passar no teste.
\item \textbf{Refatoração:} Melhore a qualidade sem quebrar o teste.
\end{enumerate}

\subsection*{Exercício 3.3 (Nível Avançado) -- EventBus com Prioridade}
Analise o código do \codigo{EventBus} e proponha uma extensão para
suportar eventos com prioridade.

\textbf{Solução (esboço de implementação):}
\begin{lstlisting}[style=pythonstyle]
import heapq
from dataclasses import dataclass, field

@dataclass(order=True)
class EventoPrioritario:
    prioridade: int
    timestamp: float = field(compare=False)
    tipo: str = field(compare=False)
    dados: dict = field(compare=False)

class EventBusComPrioridade:
    def __init__(self):
        self._heap = []
        self._handlers = {}

    def publish(self, tipo, dados, prioridade=0):
        evento = EventoPrioritario(prioridade, time.time(), tipo, dados)
        heapq.heappush(self._heap, evento)

    def process(self):
        while self._heap:
            evento = heapq.heappop(self._heap)
            if evento.tipo in self._handlers:
                self._handlers[evento.tipo](evento.dados)
\end{lstlisting}

\subsection*{Exercício 3.4 (Nível PhD) -- Novo Tipo de Raciocínio}
Implemente um novo tipo de raciocínio (raciocínio probabilístico) e
registre-o no catálogo do Nexus.

\textbf{Solução (template conceitual):}
\begin{lstlisting}[style=pythonstyle]
class ProbabilisticReasoning:
    tipo = "probabilistico"
    categoria = "decisao"

    def raciocinar(self, premissas: dict) -> dict:
        """Usa teorema de Bayes para inferencia."""
        prob_prior = premissas.get("prior", 0.5)
        verossimilhanca = premissas.get("likelihood", 1.0)
        evidencia = premissas.get("evidence", 1.0)
        prob_posterior = (verossimilhanca * prob_prior) / evidencia
        return {"conclusao": prob_posterior, "tipo": self.tipo}
\end{lstlisting}

Registre no catálogo: \codigo{nexus.catalog.register(ProbabilisticReasoning())}.

% =============================================================================
%  CAPÍTULO 4: SCANNER E METACOGNIÇÃO
% =============================================================================
\section*{Capítulo 4 -- Scanner Pipeline e Metacognição}
\addcontentsline{toc}{section}{Capítulo 4}

\subsection*{Exercício 4.1 (Nível 0) -- A Analogia do ``Olho que se Vê''}
Explique, com suas próprias palavras, a analogia do ``olho que se vê''.
Por que um sistema de software precisa de auto-observação?

\textbf{Solução:} A analogia remete ao paradoxo de um olho que tenta
observar a si mesmo: o ato de observar já modifica o observado.
No contexto de software, um sistema que se auto-observa pode detectar
seus próprios padrões, identificar gargalos, pontos cegos e vícios
epistemológicos. Um sistema sem auto-observação é como um navio sem
instrumentos: navega, mas não sabe se está no rumo certo. A
auto-observação (metacognição) é essencial para auto-evolução dirigida.

\subsection*{Exercício 4.2 (Nível Intermediário) -- Teleological Coverage Score}
Implemente o \emph{Teleological Coverage Score} (TCS).

\textbf{Solução:}
\begin{lstlisting}[style=pythonstyle]
def teleological_coverage_score(
    categorias_requeridas: dict[str, float],
    categorias_presentes: set[str]
) -> float:
    """TCS = soma ponderada das categorias presentes / soma total dos pesos."""
    total = sum(categorias_requeridas.values())
    if total == 0:
        return 1.0
    coberto = sum(peso for cat, peso in categorias_requeridas.items()
                  if cat in categorias_presentes)
    return coberto / total

# Exemplo
requeridas = {"diagnostico": 0.5, "intervencao": 0.3, "monitoramento": 0.2}
presentes = {"diagnostico", "monitoramento"}
tcs = teleological_coverage_score(requeridas, presentes)
print(f"TCS = {tcs:.2f}")  # 0.70 (70% de cobertura teleologica)
\end{lstlisting}

\subsection*{Exercício 4.3 (Nível PhD) -- Pipeline Completo}
Execute o pipeline completo de scanners e responda aos itens do enunciado.

\textbf{Solução (roteiro de execução):}
\begin{enumerate}
\item \textbf{Noological}: \codigo{/scan noological} --- identifica pontos cegos
      nas 10 dimensões epistemológicas.
\item \textbf{Teleological}: \codigo{/scan teleological --goals goals.json}
      --- mapeia lacunas entre estado atual e desejado.
\item \textbf{Evolutionary}: \codigo{/scan evolutionary} --- gera roadmap
      com quick wins, foundations, frontiers e convergents.
\item \textbf{Refinement + MCSP}: \codigo{/scan refine} --- encontra o
      conjunto mínimo de capacidades a adquirir.
\end{enumerate}

Ao final, \codigo{/evolve} executa o roadmap e reexecuta o Noological
Scanner para verificar a redução de lacunas.

\subsection*{Exercício 4.4 (Nível PhD) -- Auto-Evolução Dialética}
Use o DialecticalEngine para identificar uma limitação no pipeline.

\textbf{Solução:} Exemplo de ciclo dialético aplicado ao pipeline:

\begin{itemize}
\item \textbf{Tese}: ``O pipeline detecta gaps epistemológicos com precisão.''
\item \textbf{Antítese}: ``O pipeline não detecta gaps de desempenho
computacional (latência, consumo de memória).''
\item \textbf{Síntese (aufheben)}: ``Estender o pipeline com um scanner
de desempenho que, a cada execução, meça latência e memória de cada
módulo e alimente esses dados como nova dimensão para o Noological
Scanner.''
\end{itemize}

Implementação: criar \codigo{PerformanceScanner} como módulo M6 que
herda da interface \codigo{BaseScanner} e se insere antes do MCSP.

% =============================================================================
%  CAPÍTULO 5: TRUST ENGINE E GOVERNANÇA
% =============================================================================
\section*{Capítulo 5 -- Trust Engine e Governança Comportamental}
\addcontentsline{toc}{section}{Capítulo 5}

Os exercícios do Capítulo~5 envolvem a implementação e validação dos
componentes do Trust Engine (SPEC-038). Abaixo, modelos de resposta
para os principais tipos de exercício.

\subsection*{Exercício 5.1 -- TrustScorer (Nível Intermediário)}
Implemente o TrustScorer com blend 70/30 entre resultado imediato e
histórico.

\textbf{Solução:}
\begin{lstlisting}[style=pythonstyle]
from dataclasses import dataclass, field

@dataclass
class TrustScorer:
    alpha: float = 0.7

    def score(self, outcome: float, history: list[float]) -> float:
        if not history:
            return outcome
        media_historica = sum(history) / len(history)
        return self.alpha * outcome + (1 - self.alpha) * media_historica

# Teste
ts = TrustScorer()
s = ts.score(0.85, [0.70, 0.72, 0.68])
print(f"Trust score: {s:.2f}")  # 0.79
\end{lstlisting}

\subsection*{Exercício 5.2 -- NaturalForgetting (Nível Avançado)}
Modele o esquecimento natural segundo Atkinson-Shiffrin.

\textbf{Solução:} O modelo de Atkinson-Shiffrin tem três estágios:
sensorial $\to$ curto prazo $\to$ longo prazo. No \opencodeeco{}:
\begin{lstlisting}[style=pythonstyle]
@dataclass
class NaturalForgetting:
    capacidade_cp: int = 7
    taxa_esquecimento: float = 0.1

    def transferir(self, item, repeticoes: int) -> bool:
        return repeticoes >= 3  # consolida na memoria de longo prazo

    def recuperar(self, item, tempo_desde_ultimo_acesso: float) -> float:
        return exp(-self.taxa_esquecimento * tempo_desde_ultimo_acesso)
\end{lstlisting}

% =============================================================================
%  CAPÍTULO 6: TOKEN ECONOMY
% =============================================================================
\section*{Capítulo 6 -- Token Economy e Sustentabilidade Econômica}
\addcontentsline{toc}{section}{Capítulo 6}

Os exercícios do Capítulo~6 exigem implementação dos componentes da
Token Economy (SPEC-022 a SPEC-024). Seguem os modelos de solução.

\subsection*{Exercício 6.1 -- Ledger (Nível Intermediário)}
Implemente o \codigo{Ledger} congelado (frozen dataclass).

\textbf{Solução:}
\begin{lstlisting}[style=pythonstyle]
from dataclasses import dataclass

@dataclass(frozen=True)
class Transacao:
    origem: str
    destino: str
    valor: int
    timestamp: float
    hash_anterior: str

    def hash(self) -> str:
        import hashlib
        conteudo = f"{self.origem}{self.destino}{self.valor}"
        return hashlib.sha256(conteudo.encode()).hexdigest()
\end{lstlisting}

\subsection*{Exercício 6.2 -- Staking (Nível Avançado)}
Implemente o staking com lock de 7 dias e slashing.

\textbf{Solução:}
\begin{lstlisting}[style=pythonstyle]
from datetime import datetime, timedelta

class StakingManager:
    def __init__(self):
        self.stakes = {}  # agente -> (valor, data_inicio)
        self.PERIODO_LOCK = timedelta(days=7)
        self.PENALIDADE_SLASH = 0.3

    def staking(self, agente: str, valor: int) -> bool:
        self.stakes[agente] = (valor, datetime.now())
        return True

    def unstaking(self, agente: str) -> int:
        valor, inicio = self.stakes.get(agente, (0, datetime.min))
        if datetime.now() - inicio < self.PERIODO_LOCK:
            penalidade = int(valor * self.PENALIDADE_SLASH)
            valor -= penalidade
        del self.stakes[agente]
        return max(valor, 0)
\end{lstlisting}

% =============================================================================
%  CAPÍTULO 7: EXPERIMENTAÇÃO E VALIDAÇÃO
% =============================================================================
\section*{Capítulo 7 -- Experimentação, Validação Científica e Produção}
\addcontentsline{toc}{section}{Capítulo 7}

Os exercícios do Capítulo~7 envolvem o CORA-Eval, Aletheia, MASWOS v5
e o PhD Auditor. Seguem os modelos de resposta.

\subsection*{Exercício 7.1 -- CORA-Eval (Nível PhD)}
Execute 150 tarefas do CORA-Eval e analise o CORA-Score.

\textbf{Solução (roteiro):}
\begin{lstlisting}[style=pythonstyle]
from cora_benchmark_tracker import CORATracker

tracker = CORATracker()
resultados = tracker.executar_todas()
cora_score = tracker.calcular_cora_score()
cora_v_score = tracker.calcular_cora_v_score(verificadores=7)
print(f"CORA-Score: {cora_score:.3f}")
print(f"CORA-V-Score: {cora_v_score:.3f}")

# Analise por dimensao
for dim, score in resultados["dimensoes"].items():
    if score < 0.5:
        print(f"ATENCAO: Dimensao {dim} com score {score:.2f}")
\end{lstlisting}

\subsection*{Exercício 7.2 -- Iterative Correction Loop (Nível Avançado)}
Execute o ciclo de correção iterativa até score $\geq 95$.

\textbf{Solução:} O ciclo segue:
\begin{enumerate}
\item Gerar artigo via MASWOS v5.
\item Avaliar com \codigo{AUTO\_SCORE\_QUALIS.py} (10 critérios).
\item Se score $< 95$, executar: 5 revisores $\to$ 4 advisors $\to$ 6
corretores linguísticos.
\item Reavaliar. Repetir até score $\geq 95$.
\item Aplicar \codigo{ptbr\_corrector.py} para remover CJK.
\end{enumerate}

% =============================================================================
%  CAPÍTULO 8: DISSERTAÇÃO E BANCA
% =============================================================================
\section*{Capítulo 8 -- Dissertação, Qualis A1 e Defesa}
\addcontentsline{toc}{section}{Capítulo 8}

Os exercícios do Capítulo~8 simulam o percurso completo de produção
acadêmica. Abaixo, guias de resposta e reflexão.

\subsection*{Exercício 8.1 -- Protocolo de Anonimato (Nível Avançado)}
Aplique o protocolo de anonimato a um anteprojeto.

\textbf{Solução (checklist):}
\begin{enumerate}
\item \textbf{Nomes próprios}: substituir por ``[Pesquisador A]''.
\item \textbf{Instituição}: substituir por ``[Instituição]''.
\item \textbf{ORCID/Lattes}: remover IDs.
\item \textbf{Orientador}: substituir por ``[Orientador]''.
\item \textbf{Afiliações de coautores}: generalizar.
\item \textbf{Localização geográfica}: remover referências a cidade/estado
       específicos (exceto quando essencial para o estudo).
\item \textbf{Agradecimentos}: remover ou generalizar.
\end{enumerate}

\subsection*{Exercício 8.2 -- Simulação de Banca (Nível PhD)}
Simule uma banca com Agent Forum e analise as perguntas geradas.

\textbf{Solução:} Configure o Agent Forum com 3 personas:
\begin{itemize}
\item \textbf{Arguidora 1} (Epistemológica): foca em fundamentação teórica
e consistência metodológica.
\item \textbf{Arguidor 2} (Técnica): questiona implementação, métricas,
reprodutibilidade.
\item \textbf{Arguidora 3} (Ética): avalia impacto social, privacidade,
vieses.
\end{itemize}

Execute: \codigo{python agent-forum/forum.py --topic dissertacao --agents 3}.

Analise as perguntas geradas. Exemplo de saída esperada:
\begin{verbatim}
[Arguidora 1] "Como voce garante que o referencial teorico adotado
nao introduz um vies de confirmacao na sua analise?"
[Arguidor 2] "Qual o intervalo de confianca dos resultados
apresentados na Tabela 4?"
[Arguidora 3] "Que salvaguardas eticas foram implementadas para
proteger os dados dos participantes?"
\end{verbatim}

\subsection*{Exercício 8.3 -- PhD Auditor (Nível PhD)}
Utilize o PhD Auditor para avaliar um artigo com quatro critérios.

\textbf{Solução (configuração e execução):}
\begin{lstlisting}[style=pythonstyle]
from mirofish.phd_auditor import PhDAuditor
from mirofish.nash_solver import NashSolver
from mirofish.statistical_rigor import CohenAnalyzer
from mirofish.qualis_a1 import QualisA1Auditor

auditor = PhDAuditor(
    nash=NashSolver(), cohen=CohenAnalyzer(),
    bonferroni=True, qualis=QualisA1Auditor()
)

resultado = auditor.avaliar(artigo_path="meu_artigo.pdf")
print(f"Score Qualis A1: {resultado.qualis_score}/100")
print(f"Equilibrio Nash: {resultado.nash_eq}")
print(f"Cohen Kappa: {resultado.cohen_kappa:.2f}")
print(f"Dimensoes significativas (Bonferroni): {
    resultado.significantes}")
\end{lstlisting}

\subsection*{Reflexão Final -- Roteiro do Nível Zero ao PhD}
O roteiro completo percorre:
\begin{enumerate}
\item \textbf{Nível 0:} Compreender os fundamentos (lógica, conjuntos).
\item \textbf{Nível Básico:} Implementar agentes reativos e regressão linear.
\item \textbf{Nível Intermediário:} Redes neurais, transformers,
      SDD+TDD, scanners.
\item \textbf{Nível Avançado:} Multiagentes, Trust Engine, Token Economy,
      CORA-Eval.
\item \textbf{Nível PhD:} Metacognição, auto-evolução, banca simulada,
      Qualis A1.
\end{enumerate}

Cada nível se apoia no anterior. A chave é a prática constante:
implementar, testar, refletir e evoluir.
